\begin{frame}{베이스라인: 무작위 정책의 리턴을 먼저 재자}
  \begin{itemize}
    \item M1 제출물의 핵심 숫자 — ``배운 정책이 실제로 \textbf{무작위보다
      나은가}''를 판정할 \kb{제로점}이 없으면 ``개선됐다''고 말할 수 없다.
    \item Pendulum-v1 매 스텝 무작위 토크 정책(100 에피소드)의 평균 리턴:
      \textbf{$-1274.1$ (sd 300.0)}, 대략 $[-1800,\,-760]$ 구간에 분포.
    \item \textbf{읽기 1} — ``0''과 ``좋음''의 거리가 표 기반보다 훨씬 크다.
      CliffWalking은 무작위 $-5000 \to$ 학습 후 $-13$ (400배 격차)였는데,
      여기선 ``좋음''이 0이고 무작위가 이미 $-1274$ — \textbf{0까지의 거리
      자체가 지표}.
    \item \textbf{읽기 2} — 아무것도 하지 않는 정책이 생각보다 나쁘지 않다.
      토크 0 고정(zero-torque)이 약 \textbf{$-1230$}으로 무작위와 같은 등급.
    \item ``학습 곡선이 $-1200$ 부근에서 시작한다''는 사실은 \textbf{학습의
      증거가 전혀 안 된다} — 이 두 숫자(무작위 $-1274$, zero-torque $-1230$)를
      먼저 재야 ``PPO가 $-673$(무작위의 절반)''이라는 판단이 성립한다.
  \end{itemize}
\end{frame}
